In a one-layer transformer on $(a+b)\bmod 23$, an operand-size curriculum did not reach 95\% held-out accuracy faster than uniform sampling at the main cell (3955 vs.\ 2945 mean censored steps; 1 vs.\ 3 of 5 seeds), and the anti-curriculum was no better. The difference is not significant with 5 seeds, and two further seeds run afterwards reversed its direction (pooled 7 seeds: \rhval{pooled7/operand-size-curriculum/wd1_f0.5/steps_to_95/mean} vs.\ \rhval{pooled7/uniform-sampling/wd1_f0.5/steps_to_95/mean} steps, 2 vs.\ 3 of 7 reaching 95\%, not significant), so no ordering showed a benefit and we do not establish a loss. What the pacing sweep does support is an effect within the curriculum family ($T_c=250$ against $T_c=1000$), which we hypothesise, without an isolating run, comes from delayed full exposure of the training set rather than from operand size. Weight decay and training fraction dominated: no run generalised at $\lambda=0$ or $f=0.4$ within 4000 steps, and the curriculum-versus-uniform sign flipped between the two cells where both always succeeded, within noise. A full study would add more moduli, longer budgets with smaller learning rates, more seeds, and other difficulty measures.
